\section{Directional coverage and small cuts}\label{sec:probability}

Lemma~\ref{lem:primary} requires frequent occurrences on the micro
layers and, from each blank term, a path along which all occurrences lie
at a single other site. To see how the two requirements fit together,
assign each micro term a color and associate it with one blank term on
each shore of its site. These two blank terms are its \emph{keys}. An
outcome selects a color and opens certain keys in that color. An
occurrence requires a term of the selected color with both keys open.

Follow the outward spine from a blank term. If it encounters an
occurrence, then at that site it also encounters an open key in its
direction of travel. It suffices, from the first such key, to find a
continuation to a leaf whose moves charge no open key outside that
site. The rules below encode this condition in a graph that compresses
each straight run of parent steps into one move. A move \emph{charges}
its direction key if the run could contain a term of the selected color.
Thus every occurrence encountered on the run forces a charge to an
open key in the occurrence's site.

The proof builds the rules on a nested sequence of intervals. Within
one interval and for a fixed color, a collection of successful keys
\emph{covers} a test if it is charged on every short path towards one
end of the interval. Low probability of coverage supplies the desired
continuations. This section bounds the weight of the small sets of keys
that can cover such a test and orders typical pairs of those sets. In
Section~\ref{sec:switching} that order will allow the preferred direction
to change as the interval grows. Section~\ref{sec:finite} realizes the
rules on terms using a finite probability space.

\subsection{Moves, colors and path tests}

The sites are the integers from $1$ to $L_0$, in their usual order.
We work in a directed acyclic graph of bounded outdegree. The graph may
be infinite. Write $m(v)$ for the site of a node $v$. A node is either
a \emph{micro node} or a \emph{blank node} on the west or east shore of its site. The
blank nodes are also called \emph{keys}. We say that $u$ reaches $v$ if
there is a finite directed path from $u$ to $v$, allowing the empty path.
We call every node reachable from $u$ an ancestor of $u$.

A blank leaf has no moves. From any other blank node there are two
basic moves. The outward move leads to a blank on the facing shore of
the adjacent site, and the inward move leads to a micro node in the
same site. Thus a blank node with no neighboring site across its shore
is a leaf. Each micro node has a basic move in either direction, leading
to a key on the corresponding shore of its own site, its
\emph{direction key}. It may also have jumps in that direction to
other micro nodes in its site. Every such jump preserves the direction
key. A \emph{spine} follows basic moves in one geographical direction
until it reaches a leaf or leaves the interval being tested. At most
three of its moves originate at a given site, and it visits at most
one micro node there.

The graph is revealed in rounds. All micro nodes have the same color
in round zero. At each later round, a deterministic rule assigns a
shore to each old color. A micro node appends to its old color a datum
computed at its key on that shore. Once these colors have been
calculated, the round supplies flags and witness jumps subject to the
following rules.
\begin{enumerate}
\item For each micro node, direction and color there is a flag. When
the flag is true, every move from that node in that direction charges
its direction key in the specified color. All other moves carry no
charge in that color. In particular a blank node never charges a key.
\item A true flag in a refined color requires a true flag in the same
direction for its projection to the preceding round.
\item Each true flag has a witness of the specified color, either the
micro node itself or a micro node reached by one jump in that
direction. The round adds at most one such jump per node, direction
and color. All earlier moves remain available.
\end{enumerate}
All flags, witnesses and calculated colors are independent of the
randomness used for the gates. In the application to free terms, a basic
move from a micro node compresses travel through several micro layers.
Its flags record the colors encountered there, and the witness jumps
recover individual occurrences. A refinement can add paths
and can remove charges when an old path is read in the refined color.
At a charging move we may always use the basic move to the charged key
instead, with the same charge. Alternatively we may reach the witness
from the source of the charging move in at most one move.

For a cluster (an interval of sites), a \emph{west test at $v$ to depth
$e$} consists of all paths from $v$ inside the cluster to a leaf or
through their first west exit. The number of moves originating at any
one site is at most $e$, including the exiting move. We also use
depth vectors, with a separate bound for each site. East tests are
defined by reflection. A set of keys \emph{covers} a test in a color if
every tested path charges some key in the set in that color. Random
\emph{coverage} means coverage by the keys whose gates in that color
succeed. Each test of finite depth is finite, since path length is bounded
by the sum of the depths and branching is bounded.

Two observations account for the depth bounds below. A depth is
charged to the \emph{source} of a move. Thus replacing a charging jump
by the basic move to its direction key preserves both the charge and
the depth spent at each site. Also, suppose a path prefix from \(v\)
to \(u\) stays in the cluster, avoids charges to a set \(B\), and spends
at most \(b_x\) moves at each site \(x\). If \(B\) covers the west test
at \(v\) to depth vector \(f\ge b\), then \(B\) covers the west test at
\(u\) to depth \(f-b\). Prepending the prefix to a missed path at \(u\)
would give a missed path at \(v\). This applies whatever directions
the prefix follows. Enlarging a depth only enlarges the set of test
paths. Thus a covering set also covers every shallower test, and a
path witnessing failure of coverage persists in every deeper test.

A rule at a key is \emph{local} if it reads a bounded outgoing unfolding
at that key, with the types of the nodes, their sites, flags, labeled
move slots and the equality relation on nodes in the unfolding. The
radius may depend on the parameters of the construction. A rule
refining a color sees only flags and moves already supplied.
The independent random bits read by a gate are indexed by keys or
bounded nonempty sets of keys in the unfolding, with auxiliary indices
such as stages and colors. Different appearances of the same index refer
to the same bit. Both the radius and the number of slots are uniformly
bounded once the parameters are fixed.

\subsection{The induction and its screens}

Fix an integer $C\ge 50$ and put
\begin{equation}\label{eq:prob-parameters}
 \begin{gathered}
 K=4C,\quad J=10000K\,4^K,\quad w=3J,\quad h=Cw,\\
 q=a^w,\quad L_0=a^h=q^C,\quad t=a^{-J},\quad U=at,\quad z=a^{-1/C}.
 \end{gathered}
\end{equation}
Here $a$ is a sufficiently large power of $2$ that is a perfect
$C$th power. Fix a constant $\lambda>0$ with
$K\exp(-\lambda/5)<1/4$. Constants in the estimates can depend on $C$
and $\lambda$. Bounds of the form $a^{o(1)}$ are uniform in the graph.

At stage \(i\), \(0\le i<h\), partition the sites from left to right
into successive clusters of \(a^i\) sites. Each cluster after stage
zero has \(a\) children from the preceding partition. Two rounds of
refinement at each positive stage assign each color a direction.
Within each cluster a key's gate uses only bits indexed in that cluster.

The induction asserts that for every node and color the probability of
coverage to per-site depth \(D_i\) in its assigned direction is at most
\(t\). Here \(D_i\ge4\) and the number of colors are polylogarithmic
in \(a\), since \(C\) and hence \(h\) are fixed. We need the same bound
after choosing any collection of \emph{ignored sites}. Fix the raw
bits at their keys arbitrarily and count only charges originating
outside those sites. All other bits retain their prescribed laws.

There is an independent raw bit at each key at each stage, shared by all
colors. At stage $0$ the gate is the raw bit, of success probability
$q^{-1}$. Assign the unique color the west direction and take $D_0=4$.
The west spine in a single site has at most one charge, so the initial
coverage probability is at most $q^{-1}\le t$, also with ignored sites.
At each positive stage a gate retains the requirements of its projected
color and adds three requirements. Its new raw bit must succeed (with
probability $z$), its \emph{screen} must pass, and a specified collection
of independent \emph{barrier bits} must all succeed. The screen uses
only randomness of previous stages. The barrier indices depend
deterministically on the key and color. The construction will give a
joint passage probability of at least $\exp(-4\lambda)$ for the
barriers at any one key and color at that stage.

Fix a positive stage $i$ and an old color (the color after stage $i-1$).
Reflect the sites, and the coordinate $m(v)$ if necessary
so that the old direction is west. Put \(d=D_{i-1}\). In a new cluster the old probability of west
coverage to depth \(d\) is at most \(U=at\), including under the
conditioning on ignored sites. Start in the child containing the
given node. Unless the test there is covered, choose a missed path
using only that child's bits. It either terminates at a leaf or exits
into the adjacent child to the west. In the second case apply the
same rule from the new starting node. Its location is determined
before any bits of the new child are inspected. Thus the conditional
probability of coverage at each step is at most \(t\). The visited
children are distinct, so concatenating the paths respects every
per-site depth. A union bound over at most \(a\) children proves
the estimate.

At a key $g$, let $A_g$ be old west coverage from $g$ in the new cluster
to depth $d$ with $g$'s own site ignored. In particular
$\Pr(A_g)\le U$ even conditional on arbitrary raw bits at its site.
A \emph{pin} at a key means that its raw bits have all succeeded through
stage $i-1$. Pins at distinct keys are independent, each with probability
\[
 p=q^{-1}z^{i-1}.
\]
For a set of keys $S$ we write $[S]$ for the event that all its pins are
present. Put $\alpha_j=U^{1/4^j}$ for $0\le j\le K$. The screen at $g$
fails if $A_g$ occurs or if
\begin{equation}\label{eq:pin-screen}
 [S]\text{ occurs for some } |S|\le K
 \text{ with }\Pr(A_g\mid [S])>\alpha_{|S|}.
\end{equation}
Passing the first part of the screen guarantees a missed path from
\(g\). Passing the second means that exposing a small set of successful
pins cannot make the absence of such paths too likely. The second test
is local as well. The event \(A_g\) depends on only finitely many local
bits. Removing pins outside that dependence leaves the conditional
probability unchanged and can only lower the threshold. Thus it
suffices to check sets \(S\) of local keys. For the rest of this
section a \emph{screened success} means that the old gate succeeds
and the current screen passes. The thinning and barrier bits of the
current stage enter in Section~\ref{sec:switching}.

\begin{lemma}[Conditional cost of a screen]\label{lem:pcost}
Suppose that a bounded set $P$ of keys satisfies
$\Pr(A_g\mid[S])\le\alpha_{|S|}$ for each $S\subseteq P$ with
$|S|\le K$. Conditional on $[P]$, the probability of
\eqref{eq:pin-screen} is $O_{|P|}(\alpha_K)$.
The same bound holds for failure of the whole screen if also
$\Pr(A_g\mid[P])\le\alpha_K$, as holds whenever $|P|\le K$.
\end{lemma}

\begin{proof}
Let $X_1,\ldots,X_N$ be independent Bernoulli variables with success
probability $p$, and $E$ an event, possibly using additional randomness,
with probability $\mu$. For $u>0$ and $u\ge2\mu$, the probability that
some set of at most $K$ present successes raises the conditional
probability of $E$ above $u$ is at most $4^{K+1}\mu/u^2$. To see this,
expand $\Pr(E\mid X_1,\ldots,X_N)$ in the orthonormal product basis
whose one-coordinate factor is $(X_j-p)/\sqrt{p(1-p)}$
\cite[Sections~8.1 and~8.4]{Odonnell2014}. The constant
coefficient is $\mu$, and the sum of squares of the coefficients is at
most $\mu$ by conditional Jensen's inequality. Conditioning a proposed
set $T$ to succeed kills all terms indexed outside $T$. Each surviving
factor is at most $p^{-1/2}$ in absolute value. If the conditional
mean exceeds $u$, some nonempty $I\subseteq T$ therefore has coefficient
of magnitude at least $2^{-K-1}u p^{|I|/2}$. Sum
$\Pr([I])=p^{|I|}$ over the sets with such a coefficient and apply the
sum-of-squares bound.

Now suppose \(T\) witnesses~\eqref{eq:pin-screen}, and put
\(S=T\cap P\), \(j=|S|\). The hypothesis rules out \(T\subseteq P\),
so \(j<K\). First condition on \([S]\). The pin indicators outside
\(P\) still form a product measure, and the conditional probability
of \(A_g\) is at most \(\alpha_j\). The outside pins of \(T\) raise
this probability above \(\alpha_{|T|}\ge\alpha_{j+1}\).
For fixed \(S\), the event that some such outside pins are present
depends only on their indicators: the conditional probabilities
used to decide which sets qualify are deterministic numbers. Hence
the further conditioning on \([P]\) leaves the probability of this
event unchanged. The product-measure estimate with
\(u=\alpha_{j+1}\ge2\alpha_j\) bounds it by
\[
 O_K\bigl(\alpha_j/\alpha_{j+1}^2\bigr)
   =O_K\bigl(\alpha_{j+1}^{\,2}\bigr)=O_K(\alpha_K).
\]
Summing over \(S\subseteq P\) proves the first assertion. The
additional hypothesis bounds the cost of \(A_g\) itself.
\end{proof}

\begin{lemma}[Joint screened successes]\label{lem:joint-gates}
If $k\le K$ distinct keys can all have screened success in the old
color, their probability of doing so is at least $\gamma p^k$, where
$\gamma=\tfrac12\exp(-8\lambda hK)$. The bound also holds if, in
addition, one independently resamples all raw bits at any one site and
requires the keys at all other sites to retain screened success.
\end{lemma}

\begin{proof}
Condition on the $k$ pins. At every required screen, joint possibility
implies that none of the deterministic conditional-probability tests on
subsets of these pins violates its threshold. For a screen of an earlier
stage this applies to the corresponding earlier pins. Later raw bits are
independent of its inputs. Lemma~\ref{lem:pcost} consequently bounds the
conditional failure probability at each screen by $O_K(\alpha_K)$,
using the same $U$ in all stages. Every barrier requires positive
values of specified bits independent of the raw bits. Overlaps in these
requirements can only increase their joint probability above the product
of their individual passage probabilities, so all pass with conditional
probability at least $\exp(-4\lambda hK)$.

For the second assertion, keep the barrier bits fixed and independently
resample all raw bits at the chosen site. The only imposed pins in this
new environment are those outside that site. Every subset of them
passed the same deterministic tests under the original possibility
assumption. The conditional failure bound therefore applies to the
additional screens as well. Their barrier requirements use exactly the
original indices and bits. Conditional on the pins, all barriers pass
with probability at least \(\exp(-4\lambda hK)\). Subtract the
conditional failure probabilities of at most \(2hK\) screens.
This subtraction requires no independence between screens and
barriers. For large \(a\) the remaining probability is at least
\(\gamma\). Multiplying by the original pin probability \(p^k\)
proves both claims.
\end{proof}

\subsection{The weight of minimal cuts}

When a set of keys covers a test, every permitted path charges at least
one of those keys. To bound the probability of coverage by successful
keys, we count small covering sets that remain minimal when the depth
is reduced.
We require the two depths because concatenating a prefix uses up part
of the depth budget.

Choose the old moves and color, or the moves and a refined color after
either round of the current stage. Keep this choice fixed when testing
coverage. In every case the success probability is measured using
the screened gates for the \emph{old} color. Given per-site depths
\(d\le e\), \(e+d\le f\), with \(\max f=a^{o(1)}\), form the
\emph{\((e,f)\) list} at \(v\) of all sets \(B\) of at most \(K\) keys
satisfying three conditions.
\begin{enumerate}
\item \(B\) covers the west test at \(v\) to depth \(f\).
\item No proper subset of \(B\) covers it to depth \(e\).
\item The keys in \(B\) can have simultaneous screened success.
\end{enumerate}
The depth inequalities are componentwise when \(e,f\) are vectors.
The west spine ensures that \(B\) is nonempty. For each \(b\in B\),
a path missed by \(B\setminus\{b\}\) to depth \(e\) must charge \(b\),
so every listed key occurs in the local unfolding. In particular the
list is finite and is determined before the bits are sampled. Assign
each set the weight \(r(B)=p^{|B|}\), the probability of its pins.
The next lemma bounds both the total weight and the weight of all
strict extensions of a fixed nonempty set \(S\).

\begin{lemma}[Mass and link of a list]\label{lem:cut-list}
For every $(e,f)$ list $\mathcal L$ and every nonempty set $S$ of keys,
\[
 \sum_{B\in\mathcal L}r(B)\le aU,
 \qquad
 \sum_{\substack{B\in\mathcal L\\ B\supsetneq S}}r(B)
       \le a^{-500K}r(S).
\]
\end{lemma}

\begin{proof}
First we bound the number of cuts on the list that can succeed in
one outcome. Fix the bits and call a key available if it has screened
success. Further keys may be excluded arbitrarily. Let \(S\) contain
at most \(K\) available keys and suppose it fails to cover to depth
\(e\). We construct at most
\begin{equation}\label{eq:cut-branch}
 K\max e+1+d
\end{equation}
candidate keys outside \(S\). Every extension of \(S\) by available
keys that covers to depth \(f\) must contain a candidate, whatever
the cardinality of the extension.

Choose a depth-$e$ path $P$ missed by $S$. Offer all available keys it
charges at sites of $S$, at most $K\max e$ in number. For extensions
using none of these keys, exclude them for the rest of this step. Take
the first charge along $P$ to any remaining available key $g$. If there
is none, no such extension covers $P$. Otherwise $g$ lies at a site
unused by $S$. Its screen guarantees an old depth-$d$ west path from
$g$ missed by the old successes outside its site. Follow $P$ up to the
charge, redirect the charging move to $g$, and append this missed path.
The result is a west test path of depth at most $e+d$ in the selected
round. The prefix before the charge avoids every still eligible extension,
and the appended old path can only lose charges under refinement.
The extension must therefore use $g$ or a key charged on the appended
path at $g$'s site. Offering these at most $1+d$ possibilities proves
\eqref{eq:cut-branch}.

Branch from the empty set until $K$ keys have been chosen or the
depth-$e$ test is covered. In the latter case minimality forces any
listed extension to equal the current set. Consequently
at most $M=(K\max e+d+2)^K=a^{o(1)}$ cuts on the list can succeed
simultaneously. For any subfamily $\mathcal F$ let $E_B$ denote either
screened success of $B$ or its joint success with one-site resampling
as in Lemma~\ref{lem:joint-gates}. Counting the successful sets gives
\begin{equation}\label{eq:cut-multiplicity}
 \gamma\sum_{B\in\mathcal F} r(B)
 \le \mathbb E\sum_{B\in\mathcal F}{\bf 1}_{E_B}
 \le M\Pr\Bigl(\bigcup_{B\in\mathcal F}E_B\Bigr).
\end{equation}
Even one successful cut implies old coverage to depth $d$, since all
old paths of that depth remain available and old charges include their
refined charges. The union probability for the full list is at most $U$,
giving the mass bound.

For the link bound suppose the family of strict extensions of $S$ on the
list is nonempty. By minimality there is a depth-$e$ path $P$ missed by
$S$. Fix one independently of the bits. The keys it charges at sites of
$S$ are all outside $S$, at most $K\max e$ in number. Conditional on
$[S]$, the probability that at least one has a pin is at most
$K(\max e)p$. For the extensions using such a key, the union of the success events
requires both $[S]$ and an extra pin on the path. By
\eqref{eq:cut-multiplicity} their weight is at most
$M K(\max e)p\,r(S)/\gamma$.

Let $\mathcal R$ be the family of remaining extensions. If it is
nonempty, let $g$ be the first key in $\bigcup_{B\in\mathcal R}B$
charged along $P$. Thus its site is free of $S$, the prefix
before its charge misses the whole family, and some member contains
$S\cup\{g\}$. Possibility for that member implies
$\Pr(A_g\mid[S])\le\alpha_{|S|}$ by the screen at $g$.
Independently resample all raw bits at $g$'s site, recomputing every
gate and screen with these replacements. Conditional on the
original pins $[S]$, the same bound on $A_g$ holds in the resampled
environment. Unless $A_g$ occurs there, choose a missed old path from
$g$ to depth $d$, using the resampled environment. For $B\in\mathcal R$
let $E_B$ denote the joint event of Lemma~\ref{lem:joint-gates},
including resampling. If $E_B$ occurs, $B$ must meet the concatenation
of the redirected prefix of $P$ with the missed path. All its keys outside
$g$'s site still succeed after resampling. Thus $B$ must contain $g$ or
one of at most $d$ keys charged on this path at $g$'s site. Conditional
on $[S]$, the original raw bits of these candidate keys are independent
of the resampled environment used to choose the path. The probability
that at least one has an original pin is at most $(1+d)p$.

Thus $\bigcup_{B\in\mathcal R}E_B$ requires $[S]$ and has conditional
probability at most $\alpha_K+(1+d)p$ given $[S]$. By
\eqref{eq:cut-multiplicity},
\[
 \gamma\sum_{B\in\mathcal R}r(B)
 \le M r(S)\bigl(\alpha_K+(1+d)p\bigr).
\]
In the two contributions the factor multiplying \(r(S)\) is at most
\[
  (M/\gamma)\bigl(\alpha_K+(K\max e+1+d)p\bigr).
\]
The factor \(M/\gamma\) and the depths are \(a^{o(1)}\).
Since \(\alpha_K=a^{-(J-1)/4^K}\), \(p\le a^{-3J}\), and
\(J=10000K\,4^K\), the displayed factor is at most \(a^{-500K}\)
for sufficiently large \(a\).
\end{proof}

Set
\begin{equation}\label{eq:cut-depths}
 H=(2K+8)d,\qquad H_-=(K+3)d,\qquad E_0=t a^{-200K}.
\end{equation}

\begin{lemma}[Coverage by a proper subset]\label{lem:proper-loss}
On any list at $v$, cuts with a proper subset covering the depth-$H_-$
west test at a specified node $u$ have total weight at most $E_0$.
The test at $u$ may use a different one of the permitted refinements
of the same old color.
\end{lemma}

\begin{proof}
For each such cut $B$, record the minimum size of a subset of $B$
covering the test at $u$ to depths $d,2d,\ldots,(K+2)d$. These numbers
are nondecreasing and lie between $1$ and $K-1$. For two consecutive
equal values, take a minimizer $S$ at the larger depth. It is minimal
even at the smaller depth, and $S\subsetneq B$ has possible simultaneous
success. Hence $S$ belongs to one of $K+1$ lists at $u$. Summing first
over $B\supsetneq S$ by the link bound and then over these lists by the
mass bound gives at most $(K+1)a^{-500K}aU\le E_0$.
\end{proof}

\subsection{The order of typical cuts}

For switching direction we need more than the weight of individual
cuts. We compare two cuts on the same list by tracing the first one
charged along a path. A \emph{core list} has \(e=H\) and \(f\ge4H\).
Call an ordered pair \(B,D\) of its cuts \emph{bad} if they share a
key, or if a proper subset of either cut covers the depth-\(H_-\)
west test from any key of the other. Use the moves and charges of the
list in these tests. Every other pair is \emph{good}.

\begin{lemma}[Good comparisons]\label{lem:good-comparisons}
The total product weight of ordered bad pairs on a core list is at most
$t^2 a^{-100K}$. For a good pair one cut, say $B$, is newer in the
following sense. Every key of $B$ reaches every key of $D$, and at every
key of $B$ the set $D$ covers the west test to depth $f-H$, minimally
even to depth $H$.
\end{lemma}

\begin{proof}
The weight of cuts containing a fixed key is at most
$(1+a^{-500K})p\le2p$ by Lemma~\ref{lem:cut-list}, including the
possible singleton. Summing over the at most $K$ keys of a cut, and using
Lemma~\ref{lem:proper-loss} in both orders, bounds the bad product weight
by $aU\,(2Kp+2K E_0)\le t^2a^{-100K}$.

Consider any path prefix within the cluster to depth $H$ which charges
one of the two cuts before it charges the other. Suppose the first charge
is to $b\in B$. Redirect that move to $b$, preserving the prefix depth
and charge. Any west test path to depth $f-H$ from $b$ can be appended
to this prefix. The resulting path at $v$ must charge $D$, and the
prefix does not. Hence $D$ covers the test at $b$. Goodness rules out
coverage by a proper subset even at depth $H_-$, hence also at depth $H$.
For each $g\in D$ there is consequently a depth-$H$ path at $b$ missed
by $D\setminus\{g\}$. That path must charge $g$. Redirecting that charge shows
that $b$ reaches $g$.

There is such a first charge on the west spine from $v$. Name the cuts
so that it belongs to $B$. No depth-$H$ prefix can then charge $D$
first. Otherwise the argument with the roles interchanged would produce a
key \(g\in D\) reaching every key of \(B\), including \(b\). Since
\(b\) also reaches \(g\) and the two cuts are disjoint, this would
give a directed cycle. For any $b'\in B$, minimality at $v$
provides a depth-$H$ path charging $b'$ and no other key of $B$.
Its first $B$ charge must precede every $D$ charge. Applying the prefix
argument again proves the assertion at $b'$.
\end{proof}

For a core list at an east-shore key $v$, let $l(B)$ be the least site
of a west-shore key of $B$. Let $r_v^+(B)$ be the site of the first key
of $B$ charged on the east spine, or $+\infty$ if that spine misses $B$. The west spine must
charge a west key at or west of $v$, and the east spine leaves $v$'s
site before its first possible charge. Thus $l(B)\le m(v)<r_v^+(B)$.
For a good pair with $B$ newer than $D$ we have
\begin{equation}\label{eq:cut-endpoints}
 l(D)<l(B),\qquad r_v^+(B)\le r_v^+(D),
\end{equation}
and equality on the right requires both endpoints to be infinite.
Indeed the west spine from a west key realizing $l(B)$ must charge $D$
and can do so only on a strictly more westerly site. On the east spine
from $v$, the prefix argument of Lemma~\ref{lem:good-comparisons}
excludes a first $D$ charge before $B$. The spine has at most one charged
key per site, and the two cuts are disjoint.
